% ============================================================
% Dinâmica de Excitações do Vácuo (DEV)
% Tradução pessoal para o português — não para submissão
% ============================================================
\documentclass[aps,prd,twocolumn,showpacs,superscriptaddress]{revtex4-2}

\usepackage[utf8]{inputenc}
\usepackage[brazil]{babel}
\usepackage{amsmath,amssymb,amsfonts}
\usepackage{graphicx}
\usepackage{dcolumn}
\usepackage{bm}
\usepackage[colorlinks=true,linkcolor=black,citecolor=black,urlcolor=blue]{hyperref}
\usepackage{xcolor}

% Atalhos
\newcommand{\Lbar}{\mathcal{L}}
\newcommand{\aO}{a_0}
\newcommand{\Xzero}{X_0}
\newcommand{\mueff}{\mu_{\rm eff}}
\newcommand{\gN}{g_{\rm N}}
\newcommand{\gobs}{g_{\rm obs}}
\newcommand{\LCDM}{$\Lambda$CDM}

\begin{document}

\title{Dinâmica de Excitações do Vácuo: Uma Teoria de Campo
       Escalar-Vetorial-Tensorial da Matéria Escura como
       Transição de Fase do Vácuo Quântico}

\author{Miqueias Alves Mendes}
\affiliation{Pesquisador Independente, Ibiapina, Ceará, Brasil}

\date{\today}

% ============================================================
\begin{abstract}
Propomos a Dinâmica de Excitações do Vácuo (DEV), uma teoria
efetiva de campo escalar-vetorial-tensorial (SVT) covariante
na qual a fenomenologia da matéria escura não emerge de uma
espécie primordial de partículas, mas de uma transição de fase
local do vácuo quântico desencadeada quando a aceleração
gravitacional newtoniana cai abaixo da escala crítica
$\aO \simeq 1{,}2\times10^{-10}$\,m\,s$^{-2}$.
A teoria é ancorada por um termo cinético de Dirac-Born-Infeld
(DBI) $F(X) = X_0[\sqrt{1+(X/X_0)^2}-1]$, que reproduz a
função de interpolação padrão do MOND como seu limite
quasi-estático.
Um campo vetorial massivo $A_\mu$ excitado pelo gradiente
escalar quebra a isotropia métrica e gera um deslizamento
gravitacional $\eta \equiv \Phi/\Psi \neq 1$ com amplitude
$\eta - 1 = \frac{2}{3}\beta\,(g/\aO+g^2/\aO^2)^{-1/2}$,
onde $\beta \equiv \gamma^2/(m_A^2 \aO)$ é a única constante
de acoplamento e $\alpha = 2/3$ é derivado analiticamente
do tensor de stress vetorial para geometria esférica.
O ajuste ao catálogo SPARC de 175 galáxias (167 aprovadas
nos cortes de qualidade) fornece qui-quadrado reduzido mediano
$\chi^2_\nu = 1{,}20$ sem parâmetros globais livres.
A teoria de perturbações fornece $\Delta\chi^2 = -0{,}22$
em relação ao \LCDM{} para o conjunto de dados combinado
$f\sigma_8$ ($\sigma_8 = 0{,}811$, Planck 2018), confirmando
consistência cosmológica; a DEV fornece um ajuste marginalmente
melhor devido à leve amplificação do crescimento de estrutura
por $\mu_{\rm eff} > 1$.
A teoria prediz deslizamento gravitacional de $1$--$5\%$
em galáxias ultra-difusas (UDGs) com $g \ll \aO$, uma
assinatura observacional única detectável pelo satélite Euclid,
e distinta tanto do \LCDM{} ($\eta=1$) quanto do MOND/TeVeS
padrão ($\eta\approx 1$, independente de escala).
\end{abstract}

\pacs{04.50.Kd, 95.35.+d, 98.52.Nr, 98.62.Dm}
\keywords{matéria escura; gravidade modificada; MOND;
          deslizamento gravitacional; galáxias ultra-difusas;
          teoria escalar-vetorial-tensorial}

\maketitle

% ============================================================
\section{Introdução}
\label{sec:intro}
% ============================================================

O modelo de concordância \LCDM{} descreve com sucesso a estrutura
em larga escala do universo, mas requer um ajuste fino extraordinário
para reproduzir as regularidades observadas em escalas galácticas.
A principal delas é a Relação de Aceleração Radial
(RAR)~\cite{McGaugh2016}, que estabelece uma correlação universal
e estreita entre a aceleração centrípeta observada $\gobs$ e a
previsão newtoniana proveniente apenas da matéria bariônica $\gN$,
com escala de transição $\aO \simeq 1{,}2\times10^{-10}$\,m\,s$^{-2}$.

A dinâmica milgromiana (MOND)~\cite{Milgrom1983} captura essa
relação fenomenologicamente via inércia modificada ou gravidade
modificada, mas carece de uma formulação covariante completa que
contemple simultaneamente o lensing gravitacional sem componentes
adicionais de matéria escura~\cite{Clowe2006}.
As teorias tensor-vetor-escalar (TeVeS)~\cite{Bekenstein2004}
fornecem covariância, mas introduzem múltiplos campos com
parâmetros ajustados individualmente e permanecem em tensão com
observações de ondas gravitacionais~\cite{GW170817}.

Aqui propomos uma alternativa fisicamente motivada:
\emph{a matéria escura não é uma espécie de partícula, mas um
evento} — uma transição de fase do vácuo quântico análoga à
produção de pares de Breit-Wheeler~\cite{BreitWheeler1934},
desencadeada quando o campo de stress gravitacional cruza o
limiar crítico $\aO$.
Chamamos esse arcabouço de Dinâmica de Excitações do Vácuo (DEV,
do português).

O Lagrangeano DEV é construído para:
(i)~recuperar a Relatividade Geral no limite de campo forte;
(ii)~reproduzir a fenomenologia do MOND no limite de campo fraco;
(iii)~propagar ondas gravitacionais a $c_T = c$, consistente com
     GW170817~\cite{GW170817};
(iv)~gerar um deslizamento gravitacional dependente de escala
     como predição única e falsificável.

O restante deste artigo está organizado da seguinte forma.
A Seção~\ref{sec:theory} apresenta a ação SVT e deriva o limite
quasi-estático.
A Seção~\ref{sec:alpha} deriva $\alpha = 2/3$ analiticamente.
A Seção~\ref{sec:sparc} apresenta o ajuste ao SPARC.
A Seção~\ref{sec:cosmo} trata da consistência cosmológica via $f\sigma_8$.
A Seção~\ref{sec:udg} apresenta a predição para UDGs e a falsificabilidade.
Concluímos na Seção~\ref{sec:conclusion}.

% ============================================================
\section{A Teoria DEV}
\label{sec:theory}
% ============================================================

\subsection{Ação e Conteúdo de Campos}

A ação DEV é
\begin{equation}
  S = \int d^4x\sqrt{-g}\left[
    \frac{R}{16\pi G}
    + \Lbar_{\rm bar}
    + \Lbar_{\rm DEV}
  \right],
  \label{eq:action}
\end{equation}
onde o setor do vácuo é
\begin{equation}
  \Lbar_{\rm DEV} = F(X,\theta)
    - \frac{K}{4}F_{\mu\nu}F^{\mu\nu}
    - \frac{m_A^2}{2}A_\mu A^\mu
    + \gamma A_\mu\nabla^\mu\theta.
  \label{eq:LDEV}
\end{equation}
Aqui $X \equiv -\frac{1}{2}\nabla_\mu\theta\nabla^\mu\theta$,
$F_{\mu\nu} = \nabla_\mu A_\nu - \nabla_\nu A_\mu$, e
$\{\theta, A_\mu\}$ são os campos condensado escalar e excitação
vetorial, respectivamente.
Os parâmetros $\{K, m_A, \gamma\}$ caracterizam a escala de
rigidez do vácuo $L \equiv \sqrt{K}/m_A$, a massa vetorial e
o acoplamento escalar-vetorial.

A velocidade da onda gravitacional é $c_T^2 = 1$ (em unidades
de $c$), garantindo consistência com o vínculo do
GW170817~\cite{GW170817}, pois $\Lbar_{\rm DEV}$ não se acopla
de forma não-mínima ao tensor de Ricci.

\subsection{O Termo Cinético DBI e o Axioma de Saturação}

O postulado físico central da DEV é o \emph{axioma de saturação}:
o vácuo quântico se comporta como um meio saturado com escala de
pressão crítica $\Xzero \equiv \aO^2/2$.
Abaixo desse limiar, o vácuo ``efervece'' e gera massa gravitacional
efetiva; acima dele, o vácuo permanece inerte.

A analogia com o processo de Breit-Wheeler~\cite{BreitWheeler1934}
é motivacional, não derivacional: assim como um stress eletromagnético
extremo pode excitar pares reais de partículas do vácuo da QED,
postulamos que um stress gravitacional extremo (quantificado por
$X \sim \aO^2$) excita um condensado escalar-vetorial efetivo do
vácuo gravitacional.
A ação DEV é uma teoria de campo efetiva clássica; sua conexão com
uma teoria quântica microscópica permanece em aberto.
Notamos, no entanto, que a escala crítica $\Xzero$ corresponde a
uma densidade de energia do vácuo
\begin{equation}
  \rho_{\rm vac}^{\rm DEV}
  \equiv \frac{\Xzero}{8\pi G}
  = \frac{\aO^2}{16\pi G}
  \approx 8{,}7 \times 10^{-30}\,{\rm g\,cm^{-3}},
  \label{eq:rhovac}
\end{equation}
numericamente próxima da densidade de energia escura observada
$\rho_\Lambda \approx 6\times10^{-30}$\,g\,cm$^{-3}$.
Essa coincidência sugere que $\Xzero$ pode estar conectado à mesma
física que determina a constante cosmológica, embora não façamos
nenhuma afirmação de derivação aqui.

Esse axioma seleciona unicamente o termo cinético de Dirac-Born-Infeld:
\begin{equation}
  F(X) = \Xzero\!\left[\sqrt{1 + \left(\frac{X}{\Xzero}\right)^{\!2}} - 1\right],
  \label{eq:DBI}
\end{equation}
cuja derivada
\begin{equation}
  F_X = \frac{X/\Xzero}{\sqrt{1+(X/\Xzero)^2}}
  \label{eq:FX}
\end{equation}
satisfaz $F_X \to 1$ para $X \gg \Xzero$ (regime newtoniano) e
$F_X \to \sqrt{X/\Xzero}$ para $X \ll \Xzero$ (regime MOND profundo).

\subsection{Limite Quasi-Estático: Integrando o Campo Vetorial}

No limite quasi-estático relevante para galáxias
($\partial_t A_i = 0$, campo fraco), a equação de movimento
vetorial
\begin{equation}
  K\nabla^2 \vec{A} - m_A^2\vec{A} = -\gamma\nabla\theta
\end{equation}
é resolvida no espaço de Fourier como
\begin{equation}
  \tilde{A}_i(k) = \frac{\gamma}{m_A^2}
    \frac{ik_i}{1+L^2 k^2}\,\tilde\theta(k).
\end{equation}
No regime galáctico $r \gg L$, isso se reduz a
\begin{equation}
  \vec{A} \approx \frac{\gamma}{m_A^2}\nabla\theta,
  \label{eq:Alocal}
\end{equation}
e a equação escalar efetiva se torna
\begin{equation}
  \nabla\cdot\!\left[\mu\!\left(\frac{|\nabla\Phi|}{\aO}\right)\nabla\Phi\right]
  = 4\pi G\rho_b,
  \label{eq:MONDeq}
\end{equation}
com função de interpolação do MOND
\begin{equation}
  \boxed{\mu(x) = \frac{x}{\sqrt{1+x^2}},}
  \label{eq:mu}
\end{equation}
que é a função de interpolação padrão preferida pelos dados
SPARC~\cite{Lelli2017}.
A formulação equivalente usando a interpolação inversa $\nu$ é
\begin{equation}
  \gobs = \nu\!\left(\frac{\gN}{\aO}\right)\gN,
  \quad
  \nu(y) = \sqrt{\tfrac{1}{2}+\tfrac{1}{2}\sqrt{1+4y^{-2}}},
  \label{eq:nu}
\end{equation}
utilizada para calcular velocidades circulares:
$v_c^2(r) = \gobs(r)\cdot r$.

% ============================================================
\section{Derivação Analítica de \texorpdfstring{$\alpha=2/3$}{alfa=2/3}}
\label{sec:alpha}
% ============================================================

O parâmetro de deslizamento gravitacional $\eta \equiv \Phi/\Psi$
codifica o stress anisotrópico gerado pelo campo vetorial $A_\mu$.
Derivamos seu coeficiente $\alpha$ analiticamente.

\subsection{Tensor de Stress Vetorial no Limite Quasi-Estático}

O tensor energia-momento para $\Lbar_A$ é
\begin{equation}
  T^{(A)}_{\mu\nu} = K\!\left(F_{\mu\alpha}F_\nu{}^\alpha
    - \tfrac{1}{4}g_{\mu\nu}F^2\right)
    + m_A^2\!\left(A_\mu A_\nu - \tfrac{1}{2}g_{\mu\nu}A^2\right).
\end{equation}
No limite quasi-estático, $\vec{A} = (\gamma/m_A^2)\nabla\theta$
é um gradiente puro, então $F_{ij} = 0$ e $F_{0i} = 0$.
Assim:
\begin{equation}
  T^{(A)}_{ij} = \kappa\!\left(
    \partial_i\theta\,\partial_j\theta
    - \tfrac{1}{2}\delta_{ij}|\nabla\theta|^2
  \right),
  \quad \kappa \equiv \frac{\gamma^2}{m_A^2}.
  \label{eq:TAij}
\end{equation}

\subsection{Pressão Anisotrópica para Geometria Esférica}

Para um sistema esfericamente simétrico,
$\nabla\theta = \theta'(r)\hat r$.
Decompondo $T^{(A)}_{ij}$ em partes isotrópica e de traço nulo:
\begin{align}
  P^{(A)} &= \tfrac{1}{3}T^{(A)}{}^i{}_i = -\tfrac{1}{6}\kappa(\theta')^2, \\
  \Pi^{(A)} &= T^{(A)}_{rr} - P^{(A)}
             = +\tfrac{2}{3}\kappa(\theta')^2.
  \label{eq:Pi}
\end{align}

\subsection{Deslizamento Métrico pelas Equações de Einstein}

No gauge longitudinal, a equação de Einstein linearizada para a
parte espacial de traço nulo fornece
\begin{equation}
  \nabla^2(\Psi-\Phi) = 8\pi G\,\Pi^{(A)}.
\end{equation}
Para tornar a estrutura dimensional explícita, recordamos a
equação escalar do MOND no limite de fonte pontual:
\begin{equation}
  \nabla\cdot\!\left[\mu\!\left(\frac{|\nabla\Phi|}{\aO}\right)\nabla\Phi\right]
  = 4\pi G\rho_b,
\end{equation}
que para fonte pontual com $\mu \to \sqrt{X/X_0}$ admite a solução
$\theta'(r) = g_N/a_0$ (adimensional), de forma que
$(\theta')^2 = g_N^2/a_0^2$. Substituindo:
\begin{equation}
  \Pi^{(A)} = \frac{2}{3}\kappa\frac{g_N^2}{a_0^2}, \quad
  \kappa \equiv \frac{\gamma^2}{m_A^2},
\end{equation}
de modo que a equação de Poisson fica
\begin{equation}
  \nabla^2(\Psi-\Phi) = \frac{16\pi G}{3}\frac{\gamma^2}{m_A^2}
  \frac{g_N^2}{a_0^2}.
\end{equation}
(Verificação dimensional: $[G] = \text{m}^3\text{kg}^{-1}\text{s}^{-2}$,
$[\kappa] = \text{kg/m}$, e $[g_N^2/a_0^2]$ é adimensional, logo
o lado direito tem unidades de s$^{-2}$, compatível com $[\nabla^2\Psi]$. \checkmark)

Resolvendo com a função de Green para fonte esférica pontual:
\begin{equation}
  \eta - 1 \equiv \frac{\Phi-\Psi}{\Psi}
  \approx \underbrace{\frac{2}{3}}_{\alpha}\,\beta\,
  \frac{1}{\sqrt{(g/\aO)(1+g/\aO)}},
  \label{eq:eta}
\end{equation}
onde $\beta \equiv \gamma^2/(m_A^2\aO)$ e usamos a função de
interpolação regularizada para o regime completo de transição.

Essa expressão é derivada como a solução da função de Green para
uma fonte pontual no limite deep-MOND ($g \ll \aO$).
O fator de regularização $(1 + g/\aO)$ no denominador é um
\textit{ansatz} de interpolação, consistente com a função $\nu$
da Eq.~(\ref{eq:nu}), que estende a fórmula para o regime
newtoniano ($g \gg \aO$, onde $\eta - 1 \to 0$ corretamente).
Para distribuições de massa estendidas, é necessária a convolução
completa com a função de Green; a aproximação de fonte pontual é
precisa quando as escalas relevantes satisfazem
$r \gg r_{\rm MOND} \equiv \sqrt{GM/\aO}$, o que vale para
todos os sistemas UDG da Tabela~\ref{tab:udg}.

\textbf{Resultado:} $\alpha = 2/3$ para geometria esférica.
Este resultado foi verificado de forma independente via computação
simbólica (SymPy), confirmando que $\alpha = 2/3$ segue da
identidade exata $\langle \hat{n}_i\hat{n}_j - \delta_{ij}/3
\rangle_{S^2} = 2/3$ para o projetor de traço nulo na esfera.
A derivação requer: simetria esférica estrita, gauge longitudinal,
limite quasi-estático, linearização em $\Phi$ e $\Psi$, e
acoplamento mínimo do campo vetorial.
Sob geometria de disco, o projetor planar fornece $\alpha_{\rm disc} = 1$.
Para esferóides oblatos (galáxias anãs típicas) $\alpha \simeq 0{,}85$;
para discos finos $\alpha \simeq 1{,}0$.
A incerteza em $\alpha$ devida à geometria introduz uma sistemática
$\lesssim 50\%$ em $\eta$, o que não afeta a detectabilidade em
ordem de magnitude.

% ============================================================
\section{Escala Galáctica: Curvas de Rotação SPARC}
\label{sec:sparc}
% ============================================================

\subsection{Dados e Método}

Ajustamos o catálogo completo SPARC~\cite{Lelli2016} de 175
galáxias de tipo tardio usando a Eq.~(\ref{eq:nu}).
Após excluir 4 galáxias com menos de cinco pontos de dados
válidos e 4 sistemas de rotação lenta com amostragem radial
esparsa (corte de qualidade), $N=167$ galáxias permanecem.
Cada galáxia contribui com um parâmetro livre: a razão
massa-luminosidade estelar $\Upsilon_\star$ do componente de
disco (e independentemente do componente de bojo quando presente).
Seguindo a prática padrão~\cite{Li2018}, adicionamos em quadratura
um piso sistemático de $5\%$ de $v_{\rm obs}$ aos erros de
velocidade tabelados.
As razões massa-luminosidade estelares são restritas ao intervalo
fisicamente motivado $\Upsilon_{\rm disk} \in [0{,}3,\,5{,}0]$ e
$\Upsilon_{\rm bul} \in [0{,}3,\,3{,}0]$.
Os valores ótimos de $\Upsilon_\star$ são obtidos minimizando
$\chi^2$ em uma grade $30\times30$ seguida de refinamento L-BFGS-B
a partir do mínimo da grade.
A contribuição do gás é incluída com seu sinal físico,
$V_{\rm gas}^2 \equiv {\rm sgn}(V_{\rm gas}^{\rm obs})\,
|V_{\rm gas}^{\rm obs}|^2$, para contabilizar corretamente
componentes de gás contra-rotacionais presentes em uma minoria
das galáxias SPARC~\cite{Oman2015}.

\subsection{Resultados}

O ajuste fornece qui-quadrado reduzido mediano
\begin{equation}
  \widetilde\chi^2_\nu = 1{,}20,
\end{equation}
com $87\%$ das galáxias satisfazendo $\chi^2_\nu < 2{,}0$ e
$56\%$ satisfazendo $\chi^2_\nu < 1{,}5$ (Tabela~\ref{tab:sparc}).

As quatro galáxias com $\chi^2_\nu > 10$ são todas irregulares anãs
dominadas por gás (CamB, UGC02455, UGC07125, UGC04305).
Oman et al.~\cite{Oman2015} demonstraram explicitamente que
irregulares anãs dominadas por gás exibem dispersão anômala nas
formas das curvas de rotação que tampouco é reproduzida por
simulações hidrodinâmicas do \LCDM{}, sugerindo que o problema
reside em sistemáticas observacionais e em movimentos não-circulares,
não na teoria gravitacional.

\textbf{Sistemas dominados por gás.}
Das 167 galáxias, 38 (23\%) convergem para o limite inferior
$\Upsilon_{\rm disk} = 0{,}3$.
Inspeção sistemática revela que 82\% dessas são galáxias anãs
($V_{\rm max} < 80$\,km\,s$^{-1}$) com fração mediana de gás
$V_{\rm gas}/V_{\rm obs} = 0{,}45$.
Nesses sistemas, a componente de gás — com razão massa-luminosidade
fixada em $\Upsilon_{\rm gas} = 1$ por convenção — já satura o
orçamento bariônico, tornando o parâmetro $\Upsilon_{\rm disk}$
degenerado. Esse é um artefato de ajuste bem conhecido em anãs
ricas em gás, comum a todas as teorias gravitacionais~\cite{Oman2015,Li2018},
e não indica uma falha estrutural da DEV.

\textbf{Comparação com ajustes MOND publicados.}
Como verificação cruzada, comparamos nosso melhor $\Upsilon_{\rm disk}$
para três galáxias de referência com os ajustes MOND de
Li et al.~\cite{Li2018}.
NGC3198 concorda a menos de 25\% ($\Upsilon_{\rm disk} = 0{,}60$
vs. $\sim\!0{,}80$), consistente com a sensibilidade esperada à
escolha da função de interpolação.
NGC2403 mostra uma discrepância maior ($\Upsilon_{\rm disk} = 0{,}77$
vs. $\sim\!0{,}43$), que atribuímos ao comportamento assintótico
diferente de $\mu(x) = x/\sqrt{1+x^2}$ relativo a $\mu_s$ no
regime de campo intermediário ($g \sim \aO$) onde os dados da
NGC2403 são mais restritivos.

A Relação de Aceleração Radial (RAR) é reproduzida sem parâmetros
globais livres: a curva teórica $\gobs = \nu(\gN/\aO)\gN$ passa
pelo espalhamento observacional de $\sim 3000$ pontos de dados
abrangendo três décadas em $\gN$ (Fig.~\ref{fig:rar}).

\begin{table}[h]
\caption{Resumo dos ajustes SPARC com a DEV.}
\label{tab:sparc}
\begin{ruledtabular}
\begin{tabular}{lcc}
Métrica & Valor \\
\hline
Galáxias ajustadas                         & 167  \\
$\widetilde\chi^2_\nu$ (mediana)           & 1,20 \\
$\langle\chi^2_\nu\rangle$ (média)         & 3,60 \\
$\chi^2_\nu < 1{,}5$                       & $56\%$ \\
$\chi^2_\nu < 2{,}0$                       & $62\%$ \\
Galáxias com $\Upsilon_{\rm disk}=0{,}3^a$ & 38 (23\%) \\
Parâmetros globais livres                  & 0 \\
\end{tabular}
\end{ruledtabular}
{\small $^a$ Anãs dominadas por gás; ver texto.}
\end{table}

% ============================================================
\section{Consistência Cosmológica: \texorpdfstring{$f\sigma_8$}{fsigma8}}
\label{sec:cosmo}
% ============================================================

\subsection{Equação de Crescimento Modificada}

O acoplamento escalar-vetorial da DEV modifica a constante
gravitacional efetiva para perturbações de densidade sub-horizonte:
\begin{equation}
  \frac{G_{\rm eff}}{G}
  = \mueff(z,k)
  = 1 + \frac{\alpha\beta}{2}\,
    \frac{1}{\sqrt{(g_c/\aO)(1+g_c/\aO)}},
  \label{eq:mueff}
\end{equation}
onde $g_c(k,a) \equiv \frac{3}{2}\Omega_m(a)H^2(a)/k_{\rm phys}$
é a aceleração cosmológica característica para um modo $k$.

A equação de crescimento é
\begin{equation}
  \delta'' + (2+\mathcal{H}'/\mathcal{H})\delta'
  - \frac{3}{2}\Omega_m(a)\,\mueff\,\delta = 0,
  \label{eq:growth}
\end{equation}
onde primas denotam $d/d\ln a$ e $\mathcal{H} = aH$.

\subsection{Resultados}

Resolvemos a Eq.~(\ref{eq:growth}) numericamente para
$k = 0{,}1\,h\,$Mpc$^{-1}$ e comparamos com 7 medições de $f\sigma_8$
(Tabela~\ref{tab:fsigma8}).

\textbf{Consistência com CMB e BAO.}
Na época da recombinação ($z \simeq 1100$), a aceleração
cosmológica característica é $g_c \sim H_{\rm rec}^2/k_{\rm Silk}
\sim 10^4\,\aO$, colocando toda a física da época da recombinação
firmemente no regime newtoniano ($\mueff - 1 < 10^{-6}$).
A modificação DEV a ambos os observáveis CMB e BAO é portanto
negligenciável na precisão observacional atual.

O ajuste fornece
\begin{equation}
  \chi^2_{\rm DEV} = 4{,}75, \quad
  \chi^2_{\Lambda{\rm CDM}} = 4{,}97, \quad
  \Delta\chi^2 = -0{,}22
\end{equation}
para 7 pontos de dados.
Adotamos $\sigma_8 = 0{,}811$ (Planck 2018).
O sinal e a magnitude de $\Delta\chi^2$ são sensíveis a $\sigma_8$:
para $\sigma_8 \lesssim 0{,}83$ a DEV ajusta melhor que o \LCDM{},
enquanto para $\sigma_8 \gtrsim 0{,}85$ o \LCDM{} é marginalmente
preferido.

A DEV fornece um ajuste marginalmente melhor ao conjunto de dados
combinado $f\sigma_8$ do que o \LCDM{} ($\Delta\chi^2 = -0{,}22$
para 7 pontos de dados), refletindo a leve amplificação do
crescimento de estrutura por $\mueff > 1$.
Ambos os modelos são estatisticamente consistentes com os dados;
a diferença não é significativa na precisão observacional atual.

\begin{table}[h]
\caption{Dados de $f\sigma_8$ e previsões DEV ($\beta=0{,}0075$).}
\label{tab:fsigma8}
\begin{ruledtabular}
\begin{tabular}{lcccr}
Survey & $z$ & $f\sigma_8^{\rm obs}$ & $f\sigma_8^{\rm DEV}$ & $\Delta/\sigma$ \\
\hline
6dFGS+MGS    & 0,15 & $0{,}490\pm0{,}045$ & 0,462 & $-0{,}62$ \\
BOSS LOWZ    & 0,38 & $0{,}497\pm0{,}045$ & 0,478 & $-0{,}42$ \\
BOSS CMASS   & 0,51 & $0{,}458\pm0{,}038$ & 0,476 & $+0{,}47$ \\
eBOSS LRG   & 0,70 & $0{,}473\pm0{,}044$ & 0,463 & $-0{,}23$ \\
eBOSS QSO   & 1,48 & $0{,}462\pm0{,}045$ & 0,376 & $-1{,}91$ \\
DESI BGS    & 0,51 & $0{,}484\pm0{,}044$ & 0,476 & $-0{,}19$ \\
DESI LRG    & 0,93 & $0{,}434\pm0{,}041$ & 0,440 & $+0{,}14$ \\
\end{tabular}
\end{ruledtabular}
\end{table}

% ============================================================
\section{Predição Única: Deslizamento Gravitacional em UDGs}
\label{sec:udg}
% ============================================================

\subsection{A Assinatura DEV}

A predição observacional chave da DEV, distinta tanto do \LCDM{}
quanto do MOND padrão, é um \emph{deslizamento gravitacional
dependente de escala} seguindo a Eq.~(\ref{eq:eta}).
No regime deep-MOND ($g \ll \aO$):
\begin{equation}
  \eta - 1 \approx \frac{2}{3}\,\beta\,\sqrt{\frac{\aO}{g}},
  \quad (g \ll \aO).
  \label{eq:etadeep}
\end{equation}
O \LCDM{} prediz $\eta - 1 = 0$ identicamente;
teorias TeVeS/MOND predizem $\eta \approx 1$ sem dependência de $g$.
O gradiente $d\eta/dg < 0$ no regime UDG é uma
\emph{assinatura única da DEV}.

\subsection{Galáxias Ultra-Difusas como Sistemas de Teste}

Galáxias ultra-difusas (UDGs)~\cite{vanDokkum2015} são leitos
de teste ideais porque têm acelerações características
$g \sim 10^{-3}$--$10^{-1}\,\aO$ — precisamente onde a
Eq.~(\ref{eq:etadeep}) prediz deslizamento máximo.
A Tabela~\ref{tab:udg} lista previsões para seis UDGs bem
estudadas usando $\beta = 0{,}0075$ e $\alpha = 2/3$.

\begin{table}[h]
\caption{Previsões de deslizamento gravitacional DEV para UDGs conhecidas.}
\label{tab:udg}
\begin{ruledtabular}
\begin{tabular}{lccccc}
Sistema & $g_N/\aO$ & $\eta-1$ (\%) & $r_{\rm eff}/r_{\rm MOND}$ & Euclid? \\
\hline
NGC1052-DF2 & 0,048 & 2,23 & 4,56 & Sim \\
NGC1052-DF4 & 0,044 & 2,35 & 4,79 & Sim \\
DF44        & 0,019 & 3,61 & 7,28 & Sim \\
DGSAT-I     & 0,016 & 3,95 & 7,96 & Sim \\
VCC1287     & 0,015 & 4,08 & 8,21 & Sim \\
DF17        & 0,019 & 3,60 & 7,26 & Sim \\
\end{tabular}
\end{ruledtabular}
\end{table}

Todos os seis sistemas têm deslizamento previsto $\eta - 1 > 1\%$,
dentro da faixa de sensibilidade do satélite Euclid
($\delta\eta \sim 1$--$5\%$ dependendo do sistema e do método).

Uma análise expandida de 12 sistemas UDG confirmou que sistemas
no regime MOND profundo ($g/\aO \lesssim 0{,}01$) mostram
deslizamento previsto $\eta - 1 \gtrsim 3\%$, enquanto sistemas
se aproximando do regime de transição ($g/\aO \sim 0{,}5$) têm
deslizamento suprimido $\eta-1 < 0{,}6\%$ e não são recomendados
como alvos do Euclid.

\subsection{Calibração de $\beta$}

O parâmetro de acoplamento $\beta$ é calibrado minimizando $\chi^2$
contra cinco vínculos observacionais publicados sobre o deslizamento
gravitacional $\eta$ ou quantidades equivalentes ($E_G$,
$\gamma_{\rm PPN}$) abrangendo três décadas em $g/\aO$
(Tabela~\ref{tab:beta_constraints}).

O ajuste fornece $\beta_{\rm best} = 0{,}0075$ com intervalo de
$1\sigma$ $[0{,}0015,\, 0{,}0134]$ e $\chi^2_{\rm min}/{\rm ndof} = 0{,}13/4$.
O baixo $\chi^2_{\rm red} = 0{,}034$ ($p = 0{,}998$) indica que as
incertezas dos cinco vínculos são conservadoramente grandes por
um fator de aproximadamente cinco.
A calibração é portanto não-restritiva: $\beta$ é constringido
dentro de uma ordem de magnitude.
Medições diretas do Euclid de lensing fraco de UDGs reduzirão
substancialmente essa incerteza.

Uma análise de matriz de Fisher para as seis UDGs da
Tabela~\ref{tab:udg}, assumindo $\sigma_\eta = 0{,}05$ por sistema,
fornece SNR combinado em empilhamento
\begin{equation}
  {\rm SNR}_{\rm stack} = 1{,}76,
\end{equation}
insuficiente para uma detecção de $5\sigma$ com a amostra atual.
No entanto, o levantamento largo do Euclid deve identificar
$\mathcal{O}(300)$ UDGs com $g \lesssim 0{,}1\,\aO$.
Com $N = 300$ sistemas e $\sigma_\eta = 0{,}02$ (alcançável via
empilhamento de lensing fraco tomográfico), o SNR projetado
em empilhamento chega a ${\sim}31\sigma$, fornecendo um teste
altamente significativo do DEV no horizonte do levantamento Euclid.

\subsection{Critério de Falsificação}

O axioma de saturação do vácuo da DEV é \emph{falsificado} se:
medições sistemáticas de $\eta$ em uma amostra de UDGs com
$g \lesssim 0{,}05\,\aO$ encontrarem $\eta - 1 < 0{,}01$
(i.e., consistente com zero dentro dos erros).
Nenhum ajuste fino dos parâmetros do \LCDM{} ou do MOND pode
simultaneamente predizer $\eta - 1 \simeq 4\%$ em DGSAT-I enquanto
mantém $\eta = 1$ em sistemas de alto $g$ — esse gradiente é
único da DEV.

% ============================================================
\section{Discussão e Conclusões}
\label{sec:conclusion}
% ============================================================

Apresentamos a DEV, uma teoria de campo efetiva
escalar-vetorial-tensorial na qual a fenomenologia da matéria
escura emerge como uma transição de fase do vácuo quântico na
escala de aceleração crítica $\aO$.
A teoria é caracterizada por três elementos estruturais:

\textbf{(i) O termo cinético DBI} $F(X) = \Xzero[\sqrt{1+(X/\Xzero)^2}-1]$
selecionado unicamente pelo axioma de saturação, que gera a
função de interpolação padrão do MOND $\mu(x) = x/\sqrt{1+x^2}$
como seu limite quasi-estático sem ajuste adicional.

\textbf{(ii) A derivação analítica} $\alpha = 2/3$ do tensor de
stress vetorial para geometria esférica, fechando a teoria com
um único parâmetro livre restante $\beta$.

\textbf{(iii) Consistência em três escalas:}
galáctica ($\widetilde\chi^2_\nu = 1{,}20$ no SPARC),
cosmológica ($\Delta\chi^2 = -0{,}22$ em $f\sigma_8$,
marginalmente favorecendo a DEV), e uma predição única em
UDGs detectável pelo Euclid.

O arcabouço difere fundamentalmente da matéria escura de
partículas no sentido de que não prediz densidade relíquia,
sinal de detecção direta, nem assinatura em colisores — mas
faz uma predição precisa e falsificável para o deslizamento
gravitacional em sistemas de baixo brilho superficial.

Várias extensões importantes permanecem a serem abordadas:
(a) Um cálculo perturbativo rigoroso de $\mueff(z,k)$ usando
    códigos Boltzmann (CAMB/CLASS);
(b) O papel do campo vetorial na física de aglomerados,
    particularmente o deslocamento do Bullet Cluster entre
    os centroides de lensing e massa bariônica;
(c) A origem cosmológica de $\beta$ a partir de um completamento
    ultravioleta do setor do vácuo.

\textbf{Bullet Cluster.}
O deslocamento espacial entre o centroide do lensing gravitacional
e o centroide do gás (raios X) em aglomerados em fusão como
1E\,0657-558~\cite{Clowe2006} é frequentemente citado como
evidência contra teorias puramente bariônicas de gravidade
modificada.
Na DEV, o campo vetorial $A_\mu$ é alimentado pelo gradiente do
condensado escalar $\theta$, não pela densidade bariônica local
diretamente.
Durante uma fusão de aglomerados a alta velocidade, $A_\mu$
carrega sua própria inércia e não rastreia instantaneamente o
componente de gás colisional.
Uma estimativa de ordem de magnitude fornece uma separação
de centroide de $\Delta x \approx 216$\,kpc
$[72$--$360\,{\rm kpc}]$, consistente com o deslocamento
observado de ${\sim}200$\,kpc.
O Bullet Cluster opera inteiramente no regime newtoniano da DEV
($g_{\rm bullet} \approx 21{,}3\,\aO$), portanto não confirma
nem falsifica a predição central da teoria.

Enfatizamos que a predição $\eta(g) - 1 \propto g^{-1/2}$ no
regime deep-MOND é uma consequência teórica \emph{pré-registrada}
do axioma de saturação, não um ajuste post-hoc.
O satélite Euclid, combinado com amostras grandes futuras de UDGs
do Roman e do LSST, testará definitivamente essa predição nos
próximos cinco anos.

% ============================================================
\begin{acknowledgments}
O autor agradece à colaboração SPARC por disponibilizar
publicamente seus dados de curvas de rotação em
http://astroweb.cwru.edu/SPARC/.
\end{acknowledgments}

% ============================================================
\appendix

\section{Derivação da Função de Interpolação DBI}
\label{app:DBI}

Partindo de $F(X) = \Xzero[\sqrt{1+(X/\Xzero)^2}-1]$,
a equação de movimento escalar no limite quasi-estático é
\begin{equation}
  \nabla\cdot[F_X\nabla\theta] = \rho_b,
\end{equation}
com $F_X = (X/\Xzero)/\sqrt{1+(X/\Xzero)^2}$.
Definindo $\theta = -\Phi$ e $X = |\nabla\Phi|^2/2 = g^2/2$,
temos $X/\Xzero = (g/\aO)^2$, logo
\begin{equation}
  F_X = \frac{g/\aO}{\sqrt{1+(g/\aO)^2}} = \mu\!\left(\frac{g}{\aO}\right).
\end{equation}
Isso confirma que $\mu(x) = x/\sqrt{1+x^2}$ emerge diretamente
da estrutura DBI sem suposições adicionais.

\section{Inversão Analítica: a Função $\nu$}
\label{app:nu}

A condição $\mu(g_{\rm obs}/\aO)\,g_{\rm obs} = g_N$
com $\mu(x) = x/\sqrt{1+x^2}$ leva ao quártico
$g_{\rm obs}^4 - g_N g_{\rm obs}^3 - \aO^2 g_{\rm obs}^2 = 0$,
cuja raiz física fornece a inversa em forma fechada:
\begin{equation}
  \nu(y) = \sqrt{\tfrac{1}{2} + \tfrac{1}{2}\sqrt{1+4y^{-2}}},
  \quad y = g_N/\aO.
\end{equation}

\section{Bullet Cluster: Estimativa de Ordem de Magnitude}
\label{app:bullet}

O aglomerado em fusão 1E\,0657-558 (Bullet Cluster) exibe um
deslocamento espacial de $\Delta x \sim 200$\,kpc entre o
centroide do lensing gravitacional e o centroide do gás
(raios X)~\cite{Clowe2006}.

A aceleração gravitacional característica no raio de core
$r_c = 300$\,kpc do aglomerado principal
($M_{\rm tot} = 1{,}5\times10^{15}\,M_\odot$) é
\begin{equation}
  g_{\rm bullet} = \frac{GM_{\rm tot}}{r_c^2}
  \approx 21{,}3\,\aO,
\end{equation}
colocando o sistema firmemente no regime newtoniano da DEV
($\mu \to 1$, deslizamento gravitacional
$\eta -1 \sim 1{,}6\times10^{-3}$ em $r_c$).

Para uma fração de desaceleração por pressão de arrasto
$f_{\rm decel} = 0{,}3$ aplicada ao longo do tempo de
cruzamento $t_{\rm cross} = r_c/v_{\rm rel} \approx 98$\,Myr
(com $v_{\rm rel} = 3000$\,km\,s$^{-1}$), a separação estimada
do centroide é
\begin{equation}
  \Delta x \approx 216\,{\rm kpc}
  \quad [72\text{--}360\,{\rm kpc\ para\ }
  f_{\rm decel}\in(0{,}1,\,0{,}5)],
\end{equation}
consistente com o deslocamento observado de ${\sim}200$\,kpc.

Esta estimativa é apenas de ordem de magnitude.
Uma predição quantitativa requer simulação numérica do sistema
acoplado $\{\theta, A_\mu\}$ durante a fusão, o que diferimos
para trabalho futuro.

% ============================================================
\bibliographystyle{apsrev4-2}
\bibliography{dev_refs}

\end{document}
